\documentclass[11pt]{article}
\usepackage[utf8]{inputenc}
\usepackage[T1]{fontenc}
\usepackage{lmodern}
\usepackage[a4paper,margin=1.05in]{geometry}
\usepackage{booktabs}
\usepackage{array}
\usepackage{parskip}
\usepackage[hidelinks]{hyperref}

\title{\bfseries What This Study Can and Cannot Tell You\\[0.35em]
\large Applying the trajectories data to a PhD-or-build decision}
\date{\normalsize 23 August 2026 \quad$\cdot$\quad revised after your corrections}

\begin{document}
\maketitle

\section*{Read this before any number below}

You collected this data, so you know the caveat, but it governs everything here:
\textbf{the median is a duration, not a probability.} It says how long the road
ran for people who reached the end. It says nothing about how many set out.

This study therefore \textbf{cannot} tell you whether a PhD or a business is
more likely to work. It never sampled anyone who tried either and failed. Both
of your options appear here only through their winners. I will not write a
sentence of the form ``the data says path A beats path B,'' because no such
sentence is supportable.

What it \emph{can} do is describe the timing and shape of careers that worked,
and on your specific question it turns out to say something quite sharp.

\section{Three corrections to my first draft}

You caught three things, and two of them changed the conclusion.

\textbf{You are 24, not 22.} Born May 2002. That shortens the runway to the
median but does not change its shape --- see below.

\textbf{The visa argument is withdrawn.} My first draft leaned on immigration
asymmetry as the tiebreaker. You closed that question on 1 August 2026: no
revenue means no exposure, and a Korean entity can take revenue lawfully if it
comes to that. Both points are sound, the ruling is recorded in
\texttt{business/notes}, and I should not have re-opened it. The argument is
removed rather than softened, because it was also weaker than I made it sound.

\textbf{You do not want to be a physics researcher.} That removes the entire
academic-prize comparison from my first draft as irrelevant to you. It matters
in a second way too: it means the PhD, if you do one, is not a career --- it is
a five-to-six-year input into something else. That reframes what it has to be
worth, and the data has an opinion about it.

\section{The runway}

\begin{table}[h]
\centering
\begin{tabular}{lrl}
\toprule
Quantity & Median & Spread \\
\midrule
Age at first hit & \textbf{33.0} & middle half 30--39, $n$=105 \\
\quad if the hit landed post-1995 & 32.0 & middle half 28.8--36.0 \\
Years from age 18 to first hit & 15.0 & middle half 12--21 \\
\textbf{Years from starting the venture that worked} & \textbf{4.2} & middle half 3.0--7.2 \\
\bottomrule
\end{tabular}
\end{table}

You are 24. The median person here who eventually crossed \$10M constant did it
at 33 --- nine years out. A quarter of them did it after 39. You are not late,
and there is no closing window. That is the first thing worth internalising,
because a decision made from urgency is worse than the same decision made
calmly, and nothing in this data supports the urgency.

\subsection{The company is the short part}

Put two of those numbers together. From 18 to the hit is a median of 15 years.
From \emph{starting the venture that worked} to the hit is 4.2.

\begin{quote}
\textbf{Roughly eleven of those fifteen years happen before the winning venture
begins.}
\end{quote}

The build is short. The long part is whatever precedes it. And because this
study dates a serial founder at their \textbf{first} company to cross rather
than their famous one, the earlier failed attempts are counted \emph{inside}
those eleven years, not excluded from them.

So the question is not ``PhD or business.'' It is \emph{what the next decade is
made of}, given that the company itself is only about four years of it.

\section{The finding that bears directly on your choice}

I built a profile table for all 235 people --- highest completed degree, field,
institution --- and cut it against commercial hits only. Prizes are excluded
entirely, so this is people who built a \$10M-constant business, which is the
thing you are considering.

\begin{table}[h]
\centering
\begin{tabular}{lrrr}
\toprule
Highest completed degree & $n$ & Median age at hit & Venture clock \\
\midrule
PhD or MD & 14 & \textbf{43.5} & 4.5 yr \\
Masters & 31 & 33.8 & 4.0 yr \\
Bachelors & 45 & \textbf{31.0} & 4.0 yr \\
High school or less & 19 & 28.5--33 & 4.0 yr \\
\bottomrule
\end{tabular}
\end{table}

Among people who built a \$10M business, the doctorate holders got there about
\textbf{twelve years later} --- and their venture clock is identical. The
doctorate does not make the company slower. It moves the starting line.

Three checks, because a gap that large invites the obvious objections:

\begin{itemize}
\item \textbf{Is it just industry?} No. Every single PhD/MD commercial winner in
this sample sits in hardware/deep tech (8) or healthcare/biotech (6).
\textbf{Zero built a software or internet business.} Restricting to those slow
industries alone, the gap persists: 43.5 against 32.5.
\item \textbf{Is it just the longer education?} No. Measured from the end of
education --- which for a PhD is already five years later --- doctorate holders
still took \textbf{15.0 years} to the hit against \textbf{8.5} for bachelors or
less. They start later \emph{and} take longer afterwards.
\item \textbf{Is the sample big enough?} No. $n$=14, and that is the honest
limit on this finding. Treat the direction as informative and the magnitude as
soft. It is also survivorship on both sides: these are the PhDs who succeeded
commercially.
\end{itemize}

Set against your situation: your ideas in \texttt{\textasciitilde/business} ---
camera arbitrage, the card tracker, the recorder bundle, the storage kit,
dropshipping --- are consumer hardware, import and distribution plays. They sit
squarely in the quadrant where the bachelors-or-less cohort hits at 31, and
nowhere near the deep-tech and biotech quadrant that is the \emph{only} place
doctorate holders show up commercially in this dataset.

\section{But switching fields is cheap, which cuts the other way}

The counterweight, and it is a real one. I split the winners by whether they
entered the domain of their eventual hit soon after finishing school, or years
later.

\begin{table}[h]
\centering
\begin{tabular}{lrrr}
\toprule
& $n$ & Median age at hit & Venture clock \\
\midrule
Entered the domain within 2 years of school & 59 & 32.5 & 4.0 yr \\
Entered it 3+ years later (career switchers) & 16 & 34.2 & 4.8 yr \\
\bottomrule
\end{tabular}
\end{table}

\textbf{One in five winners switched into their field late}, with a median
detour of 4.5 years, and it cost them under two years at the median --- a
difference the intervals cannot distinguish from zero. One person switched after
28 years and still crossed.

So: doing something unrelated for several years is \textbf{not} disqualifying,
and a physics PhD would not brand you as having missed your chance. This is the
strongest thing the data says in favour of taking the time.

What it does not say is that the detour \emph{helps}. The switchers did fine
\emph{despite} the gap, not because of it.

\section{Where that leaves the decision}

The data supports these four statements, and nothing stronger:

\begin{enumerate}
\item \textbf{You are not late.} Median age at hit is 33; you are 24; the middle
half ran to 39. No urgency is justified.
\item \textbf{The venture is about four years. The decade before it is the
variable.} Choose that decade deliberately --- it is where the time goes.
\item \textbf{For the kind of business you actually want to build, the
doctorate cohort is absent.} Zero of 14 PhD/MD commercial winners built software
or consumer distribution. Your ideas live where the 31-year-olds are.
\item \textbf{A multi-year detour is survivable.} A fifth of the winners took
one, at a cost under two years. If you want the PhD for its own sake, the data
does not say it forecloses building later.
\end{enumerate}

Which points at a fairly specific reading. \textbf{The PhD is defensible as a
thing you want, or as a hedge you value. It is not defensible as preparation for
these businesses}, because the people who built these businesses did not have
one, and the ones who did took twelve years longer to get there.

\subsection{The loop you should look at directly}

In my first draft I asked what the resume was for. You answered: \emph{the
resume is for PhD prep.}

Sit with the shape of that. The lab work exists to strengthen the PhD
application. The PhD is not something you want for itself --- you have said you
do not see yourself as a physics researcher. And on the evidence above it is not
preparation for the businesses you actually intend to build, because no
doctorate holder in this dataset built anything resembling them.

So the chain currently runs: \emph{research $\rightarrow$ application
$\rightarrow$ PhD $\rightarrow$ ?} --- and the last link is the one with nothing
written on it. That is not a criticism of the effort; the effort is real and
this study is evidence of how much of it you have. It is an observation that the
effort is currently pointed at a destination you have not endorsed.

This is worth naming plainly because it is the kind of thing that survives on
momentum. Applications have deadlines and a visible next step; a business
decision has neither, so the default outcome of not choosing is that the PhD
track wins by inertia rather than by argument. If you do the PhD, I would want
it to be because you decided to --- and the data gives you room to decide slowly,
since the median winner here was still nine years out at your age.

The one thing I would not do is spend another year on resume-building for an
application whose purpose is undecided. That year is expensive in the only
currency this study measures, and it is the one year where the alternative use
--- actually testing one of your ideas --- is cheapest.

\section{The test I would actually run}

Not a recommendation on the decision, but on how to make it. You have until
October.

The camera arbitrage and card-tracker ideas are unusually \emph{testable}: real
customers, observable margins, a short cash cycle, and no capital that you
cannot walk away from. One quarter of genuine trading --- actual units, actual
landed cost, actual buyers --- would tell you more about whether you want this
life than this entire study can, and more than another year of deliberation.

That test is also cheap in exactly the dimension the data says is scarce. It
consumes a few months of the eleven-year pre-venture window, and it resolves the
question the data cannot: not \emph{how long does it take}, but \emph{is there a
specific thing here worth eleven years.} Everyone in this dataset crossed with a
particular product for a particular buyer. ``Start a business'' is not a plan,
and the four-year venture clock does not apply to a category.

If the test says yes, the decision makes itself and the PhD becomes a cost. If
it says no, you will have learned that before committing either way, which is
worth more than either answer.

\vspace{1em}
\hrule
\vspace{0.6em}

\emph{\small One last time, because it is the misreading this dataset invites
most: nine years, or thirty-three, is how long it took the people who made it.
It is not your odds, and it is not a schedule. The people who spent the same
decade and got nothing are not in this file --- not because they were rare, but
because the study was built in a way that cannot see them.}

\end{document}
